% ═══════════════════════════════════════════════════════════════════
% XeLaTeX preamble for《AI 时代的 SRE 架构师之路》PDF
% 通过 pandoc --include-in-header 注入
% ═══════════════════════════════════════════════════════════════════

% --- 中文支持：xeCJK + ctex 风格（不走 ctexbook，保留 Pandoc book class）---
\usepackage{xeCJK}
\usepackage{fontspec}

% --- 字体候选：系统里有哪个用哪个 ---
% 正文 serif（中英混排，主衬线）
\IfFontExistsTF{Noto Serif CJK SC}{%
  \setCJKmainfont[BoldFont={Noto Sans CJK SC Bold},ItalicFont={Noto Serif CJK SC}]{Noto Serif CJK SC}%
}{%
  \IfFontExistsTF{Source Han Serif SC}{%
    \setCJKmainfont[BoldFont={Source Han Sans SC Bold}]{Source Han Serif SC}%
  }{%
    \IfFontExistsTF{Songti SC}{%
      \setCJKmainfont[BoldFont={Heiti SC}]{Songti SC}%
    }{%
      \IfFontExistsTF{PingFang SC}{%
        \setCJKmainfont{PingFang SC}%
      }{%
        % 最后 fallback —— 让构建不至于直接失败
        \setCJKmainfont{STSong}%
      }%
    }%
  }%
}

% 无衬线（标题）
\IfFontExistsTF{Noto Sans CJK SC}{%
  \setCJKsansfont[BoldFont={Noto Sans CJK SC Bold}]{Noto Sans CJK SC}%
}{%
  \IfFontExistsTF{Source Han Sans SC}{%
    \setCJKsansfont{Source Han Sans SC}%
  }{%
    \IfFontExistsTF{PingFang SC}{%
      \setCJKsansfont{PingFang SC}%
    }{%
      \setCJKsansfont{STHeiti}%
    }%
  }%
}

% 等宽（代码）
\IfFontExistsTF{JetBrains Mono}{%
  \setmonofont[Scale=0.88]{JetBrains Mono}%
}{%
  \IfFontExistsTF{Menlo}{%
    \setmonofont[Scale=0.88]{Menlo}%
  }{%
    \IfFontExistsTF{Consolas}{%
      \setmonofont[Scale=0.88]{Consolas}%
    }{%
      \setmonofont[Scale=0.88]{Courier New}%
    }%
  }%
}

% 中文等宽（可选，让代码里的中文注释不太窄）
\IfFontExistsTF{Noto Sans Mono CJK SC}{%
  \setCJKmonofont{Noto Sans Mono CJK SC}%
}{%
  \IfFontExistsTF{PingFang SC}{%
    \setCJKmonofont{PingFang SC}%
  }{}%
}

% --- 段落 / 字号 / 断行 ---
\setlength{\parskip}{0.4em}
\setlength{\parindent}{0pt}            % 西式无首行缩进，中文阅读更清晰
\linespread{1.35}
\XeTeXlinebreaklocale "zh"              % 中文断行
\XeTeXlinebreakskip = 0pt plus 1pt

% --- 标题样式（titlesec）---
\usepackage{titlesec}
\titleformat{\chapter}[display]
  {\normalfont\sffamily\huge\bfseries\color{black!85}}
  {\normalsize\sffamily\color{black!50}第 \thechapter\ 章}
  {0.5em}
  {}
\titlespacing*{\chapter}{0pt}{-8pt}{16pt}

\titleformat{\section}
  {\normalfont\sffamily\Large\bfseries\color{black!80}}
  {\thesection}{0.7em}{}
\titlespacing*{\section}{0pt}{14pt}{6pt}

\titleformat{\subsection}
  {\normalfont\sffamily\large\bfseries\color{black!75}}
  {\thesubsection}{0.6em}{}
\titlespacing*{\subsection}{0pt}{10pt}{4pt}

\titleformat{\subsubsection}
  {\normalfont\sffamily\normalsize\bfseries\color{black!70}}
  {\thesubsubsection}{0.5em}{}
\titlespacing*{\subsubsection}{0pt}{8pt}{3pt}

% --- 代码块外观 ---
\usepackage{xcolor}
\definecolor{codebg}{HTML}{F7F7F5}
\definecolor{codeframe}{HTML}{E0E0DD}

% Pandoc 用 fancyvrb + Shaded 环境；我们重定义 Shaded 给它一个浅背景
\usepackage{framed}
\usepackage{fancyvrb}
\renewenvironment{Shaded}{%
  \begin{snugshade}%
}{%
  \end{snugshade}%
}
\definecolor{shadecolor}{named}{codebg}

% 允许长代码行换行（防止横向溢出）
\fvset{%
  breaklines=true,%
  breakanywhere=true,%
  fontsize=\small,%
  commandchars=\\\{\}%
}

% --- 表格美化（longtable / booktabs）---
\usepackage{booktabs}
\usepackage{longtable}
\usepackage{array}
\usepackage{calc}
\setlength{\tabcolsep}{7pt}
\renewcommand{\arraystretch}{1.25}

% --- 链接（hyperref）---
\usepackage{hyperref}
\hypersetup{
  colorlinks=true,
  linkcolor={Maroon},
  urlcolor={NavyBlue},
  citecolor={Maroon},
  pdftitle={AI 时代的 SRE 架构师之路},
  pdfauthor={SRE to AI SRE Learning Project},
  pdfsubject={SRE + AI 架构手册},
  bookmarksopen=false,
  bookmarksnumbered=true,
  pdfstartview={Fit}
}

% --- 页眉页脚 ---
\usepackage{fancyhdr}
\pagestyle{fancy}
\fancyhf{}
\fancyhead[L]{\small\sffamily\color{black!60}AI 时代的 SRE 架构师之路}
\fancyhead[R]{\small\sffamily\color{black!60}v1.8.5}
\fancyfoot[C]{\small\thepage}
\renewcommand{\headrulewidth}{0.2pt}
\renewcommand{\footrulewidth}{0pt}

% --- 列表紧凑（避免中文项目列表过高）---
\usepackage{enumitem}
\setlist{itemsep=2pt, topsep=4pt, parsep=0pt}

% --- 引用块（admonition 未 callout 映射时 pandoc 会渲染成 blockquote）---
\usepackage{mdframed}
\renewenvironment{quote}{%
  \begin{mdframed}[%
    linewidth=2pt,%
    linecolor=black!25,%
    backgroundcolor=black!3,%
    leftline=true,%
    rightline=false,%
    topline=false,%
    bottomline=false,%
    innertopmargin=6pt,%
    innerbottommargin=6pt,%
    innerleftmargin=10pt,%
    innerrightmargin=10pt]%
    \small\itshape\color{black!70}%
}{%
  \end{mdframed}%
}

% --- 图片（防止过大溢出）---
\usepackage{graphicx}
\setkeys{Gin}{width=\linewidth,keepaspectratio}

% --- 微排版（防止标点挤在行首 / 孤行）---
\widowpenalty=10000
\clubpenalty=10000
\raggedbottom

% --- 附加：支持 emoji 回退（部分 LaTeX 系统会报 Missing character）---
% 若系统没有 emoji 字体，最好在源里少用 emoji；此处设置一个 fallback
\IfFontExistsTF{Apple Color Emoji}{%
  % XeLaTeX 不支持彩色 emoji，这里只保证不缺字符
}{}

% ═══════════════════════════════════════════════════════════════════
% END
% ═══════════════════════════════════════════════════════════════════
